% Part: second-order-logic
% Chapter: syntax-and-semantics
% Section: language-of-sol

\documentclass[../../../include/open-logic-section]{subfiles}

\begin{document}

\olfileid[pt]{sol}{syn}{lan}

\olsection{A linguagem da lógica de segunda ordem}

Assim como na lógica de primeira ordem, as expressões da lógica de
segunda ordem são construídas a partir de um vocabulário básico que
contém \emph{!!{variable}s}, \emph{!!{constant}s},
\emph{!!{predicate}s} e, às vezes, \emph{!!{function}s}.  A partir
desses elementos, junto com conectivos lógicos, quantificadores e
símbolos de pontuação como parênteses e vírgulas, formam-se
\emph{termos} e \emph{!!{formula}s}.  A diferença é que, além de
variáveis para objetos, a lógica de segunda ordem também contém
variáveis para relações e funções e permite quantificá-las. Assim, os
símbolos lógicos da lógica de segunda ordem são os da lógica de
primeira ordem, acrescidos dos seguintes:

\begin{enumerate}
\item Um conjunto !!{denumerable}s de !!{variable}s de relação de
  segunda ordem para cada aridade~$n$: $\Obj V_0^n$, $\Obj V_1^n$, $\Obj V_2^n$, \dots
\item Um conjunto !!{denumerable}s de !!{variable}s de função de segunda ordem:
  $\Obj u_0^n$, $\Obj u_1^n$, $\Obj u_2^n$, \dots
\end{enumerate}

Na lógica de primeira ordem, costuma-se incluir a !!{identity}~$\eq$.
Na lógica de primeira ordem, os símbolos não lógicos de uma
linguagem~$\Lang{L}$ são fundamentais para que possamos expressar algo
interessante. É claro que existem !!{sentences} que não usam símbolos
não lógicos, mas, apenas com~$\eq$, é difícil dizer algo interessante.
Na lógica de segunda ordem, como dispomos de uma quantidade ilimitada
de variáveis de relação e de função, podemos dizer tudo o que podemos
dizer em uma linguagem de primeira ordem, mesmo sem um estoque
especial de símbolos não lógicos.

\end{document}
